\section{Method}
\label{sec:method}

Our goal is to represent the concepts of an ontology with a small vocabulary of the ontology's own concepts, while preserving as much of their meaning as possible.
Section~\ref{sec:method:problem} states the problem.
Sections~\ref{sec:method:candidates}--\ref{sec:method:selection} select the vocabulary, Section~\ref{sec:method:tokenizer} writes each concept with it, and Section~\ref{sec:method:metrics} measures how much meaning is preserved.
Sections~\ref{sec:method:baselines}--\ref{sec:method:expert} describe the three experiments: a benchmark against graph baselines, an application to real EHR data, and an expert evaluation.

% ------------------------------------------------------------------
\subsection{Problem statement}
\label{sec:method:problem}

We model the ontology as a directed acyclic graph (DAG) $G = (V, E)$ with typed edges, $E \subseteq V \times \mathcal{R} \times V$.
An edge $(u, R, v)$ states that $v$ is a parent of $u$ ($R = \isa$) or the value of one of its attributes ($R \neq \isa$; Section~\ref{sec:snomed}).
Edges therefore point from a concept to the concepts that describe it, and acyclicity guarantees that following them never leads back to the starting concept.
The hop distance $d(u, v)$ is the number of edges on a shortest directed path from $u$ to $v$, over all relationship types.

Let $M \subseteq V$ be the concepts to represent; in our case, the concepts to which the codes of the clinical data are mapped.
Given a token budget $k$, we look for a token set $T \subseteq V$ of $k$ concepts such that every concept in $M$ can be described by the tokens it reaches, with as little loss of meaning as possible.
Because tokens are concepts of the ontology itself, every token has a known meaning that a domain expert can read.
Appendix~\ref{app:notation} lists the notation.

% ------------------------------------------------------------------
\subsection{Candidate pool}
\label{sec:method:candidates}

A concept can only be described by the concepts it reaches along the edges, and its meaning fades with each hop.
We therefore draw tokens from the concepts within $D$ hops of $M$:
\[
V_D = M \cup \{\, u \in V \setminus \{r\} : d(m, u) \le D \text{ for some } m \in M \,\},
\]
where $r$ is the root concept, which every concept reaches and which carries no meaning.
We write $G_D = (V_D, E_D)$ for the subgraph induced by $V_D$.
A larger hop radius $D$ admits more general candidates at a higher computational cost; Section~\ref{sec:setup} reports the value we use.

% ------------------------------------------------------------------
\subsection{Semantic coverage with distance decay}
\label{sec:method:score}

To select tokens, we need a score of how well a token set $T$ preserves the meaning of a concept.
Let $\mathrm{Out}(u)$ be the outgoing edges of $u$ in $G_D$ and $\deg^+(u) = |\mathrm{Out}(u)|$ its out-degree.
We define the \emph{semantic coverage} $S_h(u, T) \in [0, 1]$ of a concept $u$ recursively over the horizon $h \in \{0, \dots, D\}$, the number of hops still allowed.
At horizon~0, a concept is covered only if it is a token, $S_0(u, T) = \mathbf{1}[u \in T]$, and for $h \ge 1$,
\begin{equation}
S_h(u, T) =
\begin{cases}
1, & u \in T, \\[2pt]
0, & u \notin T \text{ and } \deg^+(u) = 0, \\[2pt]
\dfrac{\lambda}{\deg^+(u)} \displaystyle\sum_{(u,R,v) \in \mathrm{Out}(u)} S_{h-1}(v, T), & \text{otherwise,}
\end{cases}
\label{eq:coverage}
\end{equation}
where $\lambda \in [0, 1]$ is the \emph{distance decay}.
A token is fully covered, a concept that is not a token and has no outgoing edge is not covered, and any other concept takes the average coverage of its out-neighbours, multiplied by $\lambda$ for the hop.

Equivalently, $S_D(m, T)$ is the expected value of $\lambda^j$ over random walks that start at $m$ and follow a uniformly chosen outgoing edge at each step, where $j$ is the step at which the walk first meets a token (the value is 0 if it meets none within $D$ steps; Appendix~\ref{app:submodular}).
Each walk follows one strand of the concept's definition, and each hop replaces the current concept by a more general one or by one of its attribute values, so part of the meaning is lost at every hop.
With $\lambda = 1$, a walk counts fully as soon as it reaches a token, however far away; with $\lambda < 1$, near tokens count more than distant ones; with $\lambda = 0$, only concepts that are themselves tokens are covered.

The selection objective is the mean coverage of the concepts to represent,
\[
F_D(T) = \frac{1}{|M|} \sum_{m \in M} S_D(m, T).
\]
$F_D$ is normalized ($F_D(\varnothing) = 0$), monotone (adding a token never lowers it) and submodular (the gain from adding a token never grows as more tokens are selected); Appendix~\ref{app:submodular} gives the proof.

\paragraph{Status-dependent decay.}
The decay need not be the same for every hop.
In SNOMED~CT, a fully defined concept is equivalent to the conjunction of its relationships, whereas the relationships of a primitive concept are necessary but not sufficient to define it (Section~\ref{sec:snomed}).
A hop into a fully defined concept therefore rewrites the meaning without losing any of it, whereas a hop into a primitive concept is where meaning can be lost.
This is expressed by replacing $\lambda$ in Eq.~\eqref{eq:coverage} with a factor that depends on the target of the hop:
\begin{equation}
S_h(u, T) = \frac{1}{\deg^+(u)} \sum_{(u,R,v) \in \mathrm{Out}(u)} \lambda(v)\, S_{h-1}(v, T),
\qquad
\lambda(v) =
\begin{cases}
1, & v \text{ fully defined},\\
\lambda, & v \text{ primitive}.
\end{cases}
\label{eq:coverage-status}
\end{equation}
A walk then contributes the product of the factors of its hops up to its first token.
Since every factor lies in $[0, 1]$, the proof of Appendix~\ref{app:submodular} holds unchanged, and so do the selection algorithm and its guarantee (Section~\ref{sec:method:selection}).
The same holds for any factor in $[0, 1]$ attached to an edge, for instance one that depends on the relationship type.
In this paper, we use the uniform decay of Eq.~\eqref{eq:coverage}, which needs no concept status and applies to any DAG ontology.

% ------------------------------------------------------------------
\subsection{Lazy greedy token selection}
\label{sec:method:selection}

Greedy selection builds $T$ one token at a time, adding at each step the candidate $t \in V_D \setminus T$ with the largest marginal gain $\Delta(t \mid T) = F_D(T \cup \{t\}) - F_D(T)$.
Because $F_D$ is monotone and submodular, the greedy set $T_k$ is within a factor $1 - 1/e$ of the best set of the same size \cite{nemhauser1978analysis}:
\begin{equation}
F_D(T_k) \ge \Big(1 - \tfrac{1}{e}\Big) \max_{T \subseteq V_D,\ |T| \le k} F_D(T).
\label{eq:guarantee}
\end{equation}
The guarantee concerns the objective $F_D$ only; the context trees and the metrics of Section~\ref{sec:method:metrics} are computed separately, and the guarantee does not extend to them.
Submodularity also allows lazy evaluation \cite{minoux1978accelerated}.
A gain computed earlier is an upper bound on the current gain of the same candidate, so at each step only the candidate with the highest bound is re-evaluated: it is selected if its updated gain is still the highest, and is put back with its updated gain otherwise.
Lazy greedy returns the same tokens as plain greedy without re-evaluating every candidate at every step.

Each gain is computed incrementally.
Adding $t$ changes the coverage only of the concepts that reach $t$ within $D$ hops, so the change is propagated backward from $t$ along incoming edges, one hop per round, and touches only this reverse neighbourhood $R_D(t)$.
One evaluation costs $O(D \cdot |E(R_D(t))|)$, where $E(R_D(t))$ are the edges among the concepts of $R_D(t)$, instead of $O(D(|V_D| + |E_D|))$ for a full recomputation; in a sparse graph such as SNOMED~CT, $R_D(t)$ is usually small.
Appendix~\ref{app:algorithms} gives the algorithms.

Finally, greedy selection never removes a token, so the token set for any budget is a prefix of the same ranking: one run gives the vocabulary for every budget $k$.

% ------------------------------------------------------------------
\subsection{Context trees}
\label{sec:method:tokenizer}

Once $T$ is fixed, each concept of $M$ must be written with tokens.
The simplest representation is the flat set of tokens that a concept reaches, but this set loses which relationship led to which token.
In Figure~\ref{fig:context-tree}a, a disorder has two finding sites, the left kidney and the right adrenal gland.
Its flat set $\{\textit{kidney}, \textit{adrenal gland}, \textit{left}, \textit{right}\}$ cannot tell it apart from a disorder of the right kidney and the left adrenal gland.
The same loss occurs whenever a token is reached through a chain of attributes, as for the laterality of any paired body structure.

We therefore represent each concept by a \emph{context tree}, built by a forward traversal from the concept, bounded by $D$, with three rules:
\begin{itemize}
    \item \isa{} edges are \emph{transparent}: a parent describes the same aspect of the concept at a more general level, so the search continues in the current context;
    \item every other edge $(u, R, v)$ opens a \emph{subcontext} keyed by $(R, v)$, since it introduces a new aspect of the definition;
    \item each branch stops at its first token, which is attached to the current context together with its \emph{depth}, the number of hops from the concept; a branch that reaches depth $D$, or a concept without outgoing edges, without meeting a token is marked \emph{uncovered}.
\end{itemize}
In Figure~\ref{fig:context-tree}a, each finding site opens its own subcontext, so \textit{left} stays with \textit{kidney} and \textit{right} with \textit{adrenal gland}.

Algorithm~\ref{alg:expand} implements these rules; it is called once per concept as \textsc{Expand}$(m, q_0, 0)$, with an empty initial context $q_0$.
Two details matter.
The token test comes before the depth test, so a token at depth $D$ is accepted, as in Eq.~\eqref{eq:coverage}.
Attribute edges are expanded before \isa{} edges, and a pair $(R, v)$ already present in the current context is not expanded again: SNOMED~CT concepts often restate the attributes of their parents, and this rule keeps one subcontext per attribute, at the level of the child.
Figure~\ref{fig:context-tree}b traces the algorithm on a toy graph; the resulting tree, written as a nested expression, is $\big(t_A,\ R_1(t_B,\ R_2(\varnothing))\big)$, where $\varnothing$ marks an uncovered branch.

\begin{figure}[t]
\centering
\begin{minipage}[c]{0.04\linewidth}\small (a)\end{minipage}%
\begin{minipage}[c]{0.54\linewidth}
\centering
\begin{tikzpicture}[
  concept/.style={draw, rounded corners, align=center, font=\footnotesize, inner sep=3pt},
  token/.style={concept, double, double distance=1.2pt},
  rel/.style={-{Stealth[length=5pt]}},
  lab/.style={font=\scriptsize, inner sep=1.5pt}
]
\node[concept] (m) at (0,0) {$m$};
\node[concept] (lk) at (2.7,1.2) {left kidney\\structure};
\node[concept] (ra) at (2.7,-1.2) {right adrenal\\gland structure};
\node[token] (k) at (5.9,1.8) {kidney};
\node[token] (l) at (5.9,0.6) {left};
\node[token] (ag) at (5.9,-0.6) {adrenal\\gland};
\node[token] (r) at (5.9,-1.8) {right};
\draw[rel] (m) -- node[lab, sloped, above, pos=0.45] {finding site} (lk);
\draw[rel] (m) -- node[lab, sloped, above, pos=0.45] {finding site} (ra);
\draw[rel] (lk) -- node[lab, sloped, above] {\texttt{IS\_A}} (k);
\draw[rel] (lk) -- node[lab, sloped, below, pos=0.4] {laterality} (l);
\draw[rel] (ra) -- node[lab, sloped, above] {\texttt{IS\_A}} (ag);
\draw[rel] (ra) -- node[lab, sloped, below, pos=0.4] {laterality} (r);
\end{tikzpicture}
\end{minipage}\hfill
\begin{minipage}[c]{0.38\linewidth}
\begin{forest}
  for tree={font=\footnotesize, grow'=0, folder, folder indent=0.6em, inner sep=1.5pt},
  tok/.style={draw, double, double distance=1pt, rounded corners=2pt}
  [{$q_0$}
    [{finding site $[q_1]$}
      [kidney, tok]
      [{laterality $[q_2]$}
        [left, tok]]]
    [{finding site $[q_3]$}
      [{adrenal gland}, tok]
      [{laterality $[q_4]$}
        [right, tok]]]]
\end{forest}
\end{minipage}

\vspace{1.2em}
\begin{minipage}[c]{0.04\linewidth}\small (b)\end{minipage}%
\begin{minipage}[c]{0.54\linewidth}
\centering
\begin{tikzpicture}[
  concept/.style={draw, rounded corners, align=center, font=\footnotesize, inner sep=3pt, minimum width=2.2em},
  token/.style={concept, double, double distance=1.2pt},
  rel/.style={-{Stealth[length=5pt]}},
  lab/.style={font=\scriptsize, inner sep=1.5pt}
]
\node[concept] (m) at (0,0) {$m$};
\node[concept] (a) at (2.2,0.9) {$a$};
\node[concept] (b) at (2.2,-0.9) {$b$};
\node[token] (tA) at (4.6,0.9) {$t_A$};
\node[token] (tB) at (4.6,-0.35) {$t_B$};
\node[concept] (z) at (4.6,-1.45) {$z$};
\draw[rel] (m) -- node[lab, sloped, above] {\texttt{IS\_A}} (a);
\draw[rel] (m) -- node[lab, sloped, below] {$R_1$} (b);
\draw[rel] (a) -- node[lab, above] {\texttt{IS\_A}} (tA);
\draw[rel] (b) -- node[lab, sloped, above] {\texttt{IS\_A}} (tB);
\draw[rel] (b) -- node[lab, sloped, below] {$R_2$} (z);
\end{tikzpicture}
\end{minipage}\hfill
\begin{minipage}[c]{0.38\linewidth}
\begin{forest}
  for tree={font=\footnotesize, grow'=0, folder, folder indent=0.6em, inner sep=1.5pt},
  tok/.style={draw, double, double distance=1pt, rounded corners=2pt}
  [{$q_0$}
    [{$t_A$}, tok]
    [{$R_1$ $[q_1]$}
      [{$t_B$}, tok]
      [{$R_2$ $[q_2]$}
        [{uncovered ($\varnothing$)}]]]]
\end{forest}
\end{minipage}
\caption{Context trees. Left: the graph around a concept $m$; double-bordered concepts are tokens. Right: the context tree of $m$.
(a)~Why trees are needed: each \emph{finding site} edge opens its own subcontext, so \textit{left} stays with \textit{kidney} and \textit{right} with \textit{adrenal gland}; the flat set $\{\textit{kidney}, \textit{adrenal gland}, \textit{left}, \textit{right}\}$ loses this pairing.
(b)~\textsc{Expand} on a toy graph with $D = 2$ and $T = \{t_A, t_B\}$: \isa{} keeps $t_A$ in the initial context $q_0$, $R_1$ opens $q_1$, which holds $t_B$, and $R_2$ opens $q_2$, whose branch reaches depth $D$ at the non-token $z$ and is uncovered.}
\label{fig:context-tree}
\end{figure}

\begin{algorithm}[t]
\caption{\textsc{Expand}$(u, q, \ell)$: builds the context tree of a concept $m$, called as \textsc{Expand}$(m, q_0, 0)$; $\ell$ counts the hops from $m$.}
\label{alg:expand}
\begin{algorithmic}[1]
\Procedure{Expand}{$u, q, \ell$}
    \If{$u \in T$} \Comment{token test first: a token at depth $D$ is accepted}
        \State attach $u$ as a token of context $q$ at depth $\ell$
        \State \Return
    \EndIf
    \If{$\ell = D$ \textbf{or} $\mathrm{Out}(u) = \varnothing$}
        \State mark this branch of $q$ as uncovered
        \State \Return
    \EndIf
    \ForAll{$(u, R, v) \in \mathrm{Out}(u)$, attribute edges before \texttt{IS\_A} edges}
        \If{$R = \texttt{IS\_A}$}
            \State \Call{Expand}{$v, q, \ell + 1$} \Comment{transparent: same context}
        \ElsIf{$q$ has no subcontext keyed by $(R, v)$}
            \State create subcontext $q'$ of $q$ keyed by $(R, v)$
            \State \Call{Expand}{$v, q', \ell + 1$}
        \EndIf
    \EndFor
\EndProcedure
\end{algorithmic}
\end{algorithm}

% ------------------------------------------------------------------
\subsection{Evaluating a tokenization}
\label{sec:method:metrics}

A tokenizer must meet four requirements at once: represent every concept, keep its meaning, keep concepts apart, and do so with short representations.
These requirements pull against each other under a fixed budget, so no single number captures the quality of a tokenization.
We therefore evaluate every token set with four core metrics, one per requirement (Table~\ref{tab:metrics}), and compare methods on all four jointly (Section~\ref{sec:results:benchmark}).

For a concept $m$, let $n_m$ be the number of tokens in its context tree, counting a token each time it appears, and $u_m$ the number of uncovered branches.
The \emph{signature} $\sigma_m$ of the tree is its nested structure of relationships, tokens and uncovered branches, without depths or order; we write $N_s$ for the number of concepts of $M$ whose signature is $s$.
The core metrics are computed per concept and averaged over $M$.

\paragraph{Represent every concept: unknown rate.}
The unknown rate is the share of concepts without any token ($n_m = 0$).
Like out-of-vocabulary words, these concepts lose all their information, so a tokenizer should leave none.

\paragraph{Keep the meaning: proximity-coverage score.}
A represented concept still loses meaning in two ways: whole branches of its definition are lost when no token is found, and specificity is lost when the tokens found are far away.
The \emph{proximity-coverage} (PC) score combines both.
It explores the same paths as Algorithm~\ref{alg:expand}, with the same stopping rules but without merging repeated subcontexts, and weights each branch by its share of the definition: a concept splits its weight equally among its outgoing relationships.
For a concept $u$ reached after $\ell$ hops,
\begin{equation}
\phi(u, \ell) =
\begin{cases}
1 - \dfrac{\ell}{D + 1}, & u \in T, \\[4pt]
0, & u \notin T \text{ and } (\ell = D \text{ or } \deg^+(u) = 0), \\[4pt]
\dfrac{1}{\deg^+(u)} \displaystyle\sum_{(u,R,w) \in \mathrm{Out}(u)} \phi(w, \ell + 1), & \text{otherwise,}
\end{cases}
\quad
\mathrm{PC}_m = \phi(m, 0).
\label{eq:pc}
\end{equation}
Unrolling the recursion, each branch $b$ of the expansion, a path from $m$ to its first token or to an uncovered end, receives a weight $\omega_b$, the product of $1/\deg^+(u)$ over the concepts $u$ it leaves; these weights sum to 1.
With $B^{+}_m$ the branches that end at a token, at depth $\ell_b$,
\[
\mathrm{PC}_m = c_m \cdot p_m,
\qquad
c_m = \sum_{b \in B^{+}_m} \omega_b,
\qquad
p_m = \frac{1}{c_m} \sum_{b \in B^{+}_m} \omega_b \Big(1 - \frac{\ell_b}{D + 1}\Big).
\]
The \emph{branch coverage} $c_m$ is the weight of the branches that end at a token, and the \emph{proximity} $p_m$ is their weighted mean closeness.
$\mathrm{PC}_m$ equals 1 when the concept is itself a token and 0 when it reaches no token.
In Figure~\ref{fig:context-tree}b, $m$ gives half of its weight to $a$, whose branch ends at $t_A$, and half to $b$, which splits it between $t_B$ and the uncovered $z$; both tokens lie at depth 2, so $c_m = 3/4$, $p_m = 1/3$ and $\mathrm{PC}_m = 1/4$.

The weights matter when the branches of a concept fan out unevenly.
Take a concept with two relationships, one leading to ten tokens at depth 2 and the other ending without a token ($D = 9$).
An unweighted mean over its eleven branch ends gives 0.73, more than the 0.45 the concept gets when the first relationship is itself a token at depth 1, although one closer token is the better representation: the many tokens of one branch hide the missing one.
With the weights, each relationship keeps half of the score, which gives 0.40 and 0.45.
In terms of random walks, $\mathrm{PC}_m$ is the expected value of $1 - j/(D + 1)$ over the walks of Section~\ref{sec:method:score}, where $j$ is the step of the first token (0 if none), and $c_m$ equals $S_D(m, T)$ at $\lambda = 1$.
The PC score thus uses the same walks as the selection objective, with a linear rather than geometric discount, and is not independent of it.
We therefore also report the unweighted mean (Appendix~\ref{app:flat-pc}), which ranks the methods in the same way, and compare the trees with expert judgement (Section~\ref{sec:method:expert}).

\paragraph{Keep concepts apart: unique rate.}
The unique rate is the share of concepts whose signature no other concept has ($N_{\sigma_m} = 1$).
Two concepts with the same tree are indistinguishable for a model, however well each of them is covered.

\paragraph{Stay short: tokens per concept.}
The number of tokens per concept, the mean of $n_m$, measures the efficiency of the tokenizer.
It sets the length of the input sequences, and the cost of standard self-attention grows quadratically with this length \cite{vaswani2017attention}.
On EHR data, we report the corresponding cost per event (Section~\ref{sec:method:ehr}).

The four metrics cannot be optimized one at a time.
A vocabulary of generic hub concepts reaches every concept, but through long and distant branches, which lowers the PC score and lengthens the trees.
A vocabulary of very specific concepts keeps the trees short and close for the concepts it contains, but leaves others unknown.
Two concepts can both score high on PC and still share one tree.
We therefore compare methods by their average rank over the four metrics rather than on any one of them.
As secondary metrics, the \emph{uncovered branch rate} is the share of concepts with at least one uncovered branch ($u_m > 0$), which keep only part of their definition, and the \emph{vocabulary use} is the share of the $k$ selected tokens that appear in at least one context tree.
A token that no concept uses takes a slot of the vocabulary, and an embedding, without serving any concept.
Because a mean can hide a minority of badly represented concepts, we also report the distribution of the per-concept PC score and number of tokens.
\begin{table}[t]
\centering
\small
\renewcommand{\arraystretch}{1.3}
\setlength{\tabcolsep}{4pt}
\begin{tabular}{@{}llcc@{}}
\toprule
\textbf{Requirement} & \textbf{Metric} & \textbf{Per concept} & \textbf{Better} \\
\midrule
Represent every concept & Unknown rate & $\mathbf{1}[n_m = 0]$ & $\downarrow$ \\
Keep its meaning & Proximity-coverage (PC) score & $c_m \cdot p_m$, Eq.~\eqref{eq:pc} & $\uparrow$ \\
Keep concepts apart & Unique rate & $\mathbf{1}[N_{\sigma_m} = 1]$ & $\uparrow$ \\
Stay short (efficiency) & Tokens per concept & $n_m$ & $\downarrow$ \\
\midrule
Keep the whole definition (secondary) & Uncovered branch rate & $\mathbf{1}[u_m > 0]$ & $\downarrow$ \\
Use the whole vocabulary (secondary) & Vocabulary use & -- (per token set) & $\uparrow$ \\
\bottomrule
\end{tabular}
\caption{Evaluation metrics: one core metric per requirement of a tokenizer. The per-concept metrics are averaged over the concepts to represent; vocabulary use is the share of the $k$ tokens that at least one context tree contains.}
\label{tab:metrics}
\end{table}

% ------------------------------------------------------------------
\subsection{Compared selection methods}
\label{sec:method:baselines}

We compare 19 configurations (Table~\ref{tab:baselines}), which all select their tokens from the same candidate pool $V_D$.
Except for random selection, each configuration ranks $V_D$ once and takes the top $k$ candidates, so one ranking serves every budget.

\paragraph{Our method.} We run our selection with $\lambda \in \{0, 0.1, \dots, 1\}$.

\paragraph{Ablation: discrete set cover.}
The ablation uses the same lazy greedy procedure and has the same guarantee, but replaces $S_h(u, T)$ by a binary indicator: a concept counts as covered as soon as some token lies within $D$ hops of it, however far away and through however few of its branches.
This is the classical maximum coverage problem \cite{nemhauser1978analysis, feige1998threshold}.
Since the procedure is identical, the difference between the two isolates the effect of the distance-aware coverage.

\paragraph{Baselines.}
Seven baselines are well defined on a DAG: random selection, three degree-based rankings and three centrality-based rankings (Table~\ref{tab:baselines}).
Degree- and PageRank-based selection also correspond to the anchor-selection strategies of NodePiece~\cite{galkin2022nodepiece}.

\begin{table}[t]
\centering
\small
\renewcommand{\arraystretch}{1.2}
\setlength{\tabcolsep}{4pt}
\begin{tabular}{@{}>{\raggedright\arraybackslash}p{1.6cm}>{\raggedright\arraybackslash}p{3.0cm}c>{\raggedright\arraybackslash}p{7.6cm}@{}}
\toprule
\textbf{Family} & \textbf{Method} & \textbf{Configs} & \textbf{Ranks candidates by} \\
\midrule
Ours & Semantic coverage greedy & 11 & marginal gain in $F_D$, $\lambda \in \{0, 0.1, \dots, 1\}$ \\
Ablation & Discrete set cover & 1 & marginal gain in the number of concepts with a token within $D$ hops \\
Random & Random & 1 & uniform sampling; metrics averaged over 20 draws \\
\midrule
\multirow{3}{*}{Degree} & Highest degree, $d \le D$ & 1 & number of concepts that reach the candidate within $D$ hops \\
 & Highest degree, $d \le 1$ & 1 & number of concepts that reach the candidate in one hop \\
 & Most \isa{} children & 1 & number of descendants within $D$ \isa{} hops \\
\midrule
\multirow{3}{*}{Centrality} & PageRank & 1 & stationary probability of a random walk restarting uniformly on $V_D$, $\alpha = 0.85$ \cite{page1999pagerank} \\
 & Personalized PageRank & 1 & same walk, restarting uniformly on $M$ \cite{haveliwala2002topic} \\
 & Closeness centrality & 1 & inverse mean distance from the concepts that reach the candidate, with the Wasserman--Faust correction \cite{freeman1978centrality, wasserman1994social} \\
\midrule
\multicolumn{2}{@{}l}{\textbf{Total}} & 19 & \\
\bottomrule
\end{tabular}
\caption{The 19 compared configurations. All select from the candidate pool $V_D$.}
\label{tab:baselines}
\end{table}

% ------------------------------------------------------------------
\subsection{Application to EHR data}
\label{sec:method:ehr}

The benchmark treats all concepts alike, whereas in EHR data a few concepts produce most events and most concepts are rare.
We therefore apply the tokenizer to the concepts $\Mhero$ that occur in a real EHR dataset (Section~\ref{sec:setup}).
We re-run the whole pipeline on $\Mhero$ (candidate pool, then selection with $\lambda = 0.9$); the selection uses only the list of concepts, not their frequencies.
An \emph{event} is one recorded occurrence of a concept.

We compare our tokenizer with \emph{frequency truncation}, the usual frequency-based reduction of EHR vocabularies \cite{steinberg2021language}: the $k$ most frequent concepts become tokens, and every other concept is mapped to a single unknown token.
We measure:
\begin{itemize}
    \item \emph{concepts and events represented}: a concept is represented if it receives at least one token other than the unknown token, and an event is represented if its concept is;
    \item \emph{rare concepts}: concepts with fewer than 100 events, for which we report the share represented and the PC score. For truncation, the PC score follows from its definition: a kept concept is its own token, at depth 0, and scores 1, and a discarded concept has no token and scores 0;
    \item \emph{borrowed strength}: a token occurs in every event of every concept that uses it, so a rare concept is learned from the events of its tokens rather than from its own; we compare the number of events of a concept with that of its least frequent token;
    \item \emph{tokens per event}: the length of the representation of an event, counting the distinct tokens of each context and one marker per subcontext; truncation always uses one token.
\end{itemize}

% ------------------------------------------------------------------
\subsection{Expert evaluation}
\label{sec:method:expert}

The PC score measures distance in the graph, not clinical meaning.
Semanticians therefore compared, blinded, the context trees of our tokenizer ($\lambda = 0.9$, $k = 5{,}000$, selected on $\Mhero$) with those of Personalized PageRank at the same budget, using the official SNOMED~CT definition of each concept as the reference.
We chose Personalized PageRank because it is the only baseline that, like our method, targets the concepts to represent: its random walk restarts on $M$. \doubt{state the rule fixed in advance; with the four-metric ranking, Personalized PageRank is the best baseline only at low compression}

\paragraph{Concepts.}
We consider the concepts of $\Mhero$ that belong to the official SNOMED~CT release, and keep those whose two trees differ; concepts with identical trees count as ``no clear difference''.
To keep the task feasible, the two trees of a concept may hold at most 16 tokens together.
Primitive concepts are included, since leaving them out would bias the sample; raters judge them by their name and stated relationships.
A concept falls into one of three categories: kept as a token by our method and decomposed by the baseline, the reverse, or decomposed by both into different trees.
We sampled 100 concepts, stratified by category and by definition status; every rater rates the same 100 concepts, in a rater-specific order, and the side (A or B) of each method is randomized and balanced within each stratum.

\paragraph{Task and analysis.}
For each concept, raters choose ``A better'', ``no clear difference'' or ``B better'', and may flag, for each version, an incorrect or ambiguous association, absent information, lost precision (too general a token), or redundant information.
For a fully defined concept, they may also flag an official definition that does not capture its meaning.
We report the preference among the rated concepts, per stratum and weighted by stratum size, with a sign test that excludes ties and bootstrap confidence intervals; separately, we report it over all concepts of the release, where concepts with identical trees count as ties.
We also report the agreement between raters (Cohen's $\kappa$, or Fleiss' $\kappa$ for more than two raters) and between expert preference and the sign of the difference in PC score.
